\chapter{Complexity — Choice Words}
%%% generated by the blueprint pipeline from TCSlib.Complexity.TuringMachine.ChoiceWords
\section{Overview}

Declarations of the Lean module \texttt{TCSlib.Complexity.TuringMachine.ChoiceWords}.

\section{Declarations}

\begin{theorem}[Uniqueness of list concatenation split]
\lean{Turing.append_split_unique}
\label{Turing.append_split_unique}
\leanok
\difficulty{4}
Let $f:\mathbb{N}\to\mathbb{N}$ be such that $n\mapsto n+f(n)$ is strictly increasing.
If lists $x,u,y,v$ over a type $\alpha$ satisfy $|u|=f(|x|)$, $|v|=f(|y|)$, and $x$
concatenated with $u$ equals $y$ concatenated with $v$, then $x=y$ and $u=v$.
\end{theorem}

\begin{theorem}[Strict growth of a polynomial expression]
\lean{Turing.certificateLength_strictMono}
\label{Turing.certificateLength_strictMono}
\leanok
\difficulty{2}
For natural numbers $C$ and $c$, the function on natural numbers given by $n\mapsto
n+C(n+1)^c$ is strictly increasing.
\end{theorem}

\begin{theorem}[Strict increase of an exponential expression]
\lean{Turing.expCertificateLength_strictMono}
\label{Turing.expCertificateLength_strictMono}
\leanok
\difficulty{2}
For natural numbers $C$ and $c$, the function $n\mapsto n+C\,2^{(n+1)^c}$ is strictly
increasing on the natural numbers.
\end{theorem}

\begin{theorem}[A polynomial choice budget bound]
\lean{Turing.choiceBudget_le}
\label{Turing.choiceBudget_le}
\leanok
\difficulty{4}
For all natural numbers $a$, $c$, and $n$, $a(n^c+1)\le 2a(n+1)^c$.
\end{theorem}

\begin{theorem}[Run determined by the first $t$ choices]
\lean{Turing.run_take_of_halted}
\label{Turing.run_take_of_halted}
\leanok
\difficulty{4}
Let $\mathrm{run}$ map Boolean words to configurations, and let $\mathrm{halted}$ be a
predicate on configurations. Suppose that whenever a run on $w$ is halted, appending any
word $w'$ to $w$ leaves its result unchanged. Suppose also that every word of length $t$
produces a halted result. Then for every word $w$ with $t\le |w|$, its prefix of length
$t$ has length $t$, produces a halted result, and has the same run result as $w$.
\end{theorem}
